% =============================================================================
%  Vanguard + Meridian — Technical Brief v4 (Stage 6 + 7 honest edition)
%  April 19 2026 — updates for CNN learned matcher + AIS fusion
% =============================================================================
\documentclass[10pt,a4paper]{article}

\usepackage[a4paper,margin=14mm,top=12mm,bottom=14mm,footskip=8mm]{geometry}

\usepackage[T1]{fontenc}
\usepackage{lmodern}
\usepackage{amssymb}
\usepackage{helvet}
\renewcommand{\familydefault}{\sfdefault}
\usepackage{microtype}
\usepackage[none]{hyphenat}
\frenchspacing
\setlength{\parskip}{4pt}
\setlength{\parindent}{0pt}

\usepackage[table]{xcolor}
\definecolor{mCyan}{HTML}{006E8F}
\definecolor{mSafe}{HTML}{2E7D32}
\definecolor{mWarn}{HTML}{E65100}
\definecolor{mFail}{HTML}{C62828}
\definecolor{mDim}{HTML}{6A7A90}
\definecolor{mRow}{HTML}{F4F6FA}
\definecolor{mRule}{HTML}{CFD8E3}

\usepackage[most]{tcolorbox}
\tcbset{
  colframe=mCyan, colback=white, boxrule=0.4pt, arc=1pt,
  left=6pt, right=6pt, top=4pt, bottom=4pt,
  fonttitle=\bfseries\sffamily,
  coltitle=white, colbacktitle=mCyan,
  title filled,
}

\usepackage{booktabs}
\usepackage{array}
\usepackage{tabularx}
\newcolumntype{L}[1]{>{\raggedright\arraybackslash}p{#1}}

\usepackage{tikz}
\usetikzlibrary{shapes,arrows.meta,positioning,fit,backgrounds,calc}

\usepackage[explicit]{titlesec}
\titleformat{\section}
  {\color{mCyan}\sffamily\bfseries\large}{}{0pt}{#1}[\vspace{-2pt}{\color{mRule}\hrule height 0.4pt}]
\titlespacing{\section}{0pt}{10pt}{5pt}
\titleformat{\subsection}
  {\color{mCyan}\sffamily\bfseries\normalsize}{}{0pt}{#1}
\titlespacing{\subsection}{0pt}{6pt}{2pt}

\usepackage{hyperref}
\hypersetup{colorlinks,urlcolor=mCyan,linkcolor=mCyan}
\usepackage{enumitem}
\setlist[itemize]{leftmargin=*, itemsep=1pt, topsep=2pt}
\setlist[enumerate]{leftmargin=*, itemsep=1pt, topsep=2pt}

\newcommand{\statusbuilt}{{\color{white}\colorbox{mSafe}{\small\bfseries\,BUILT\,}}}
\newcommand{\statusplan}{{\color{white}\colorbox{mWarn}{\small\bfseries\,PLANNED\,}}}
\newcommand{\statusfuture}{{\color{white}\colorbox{mDim}{\small\bfseries\,FUTURE\,}}}
\newcommand{\statustuning}{{\color{white}\colorbox{mWarn}{\small\bfseries\,TUNING\,}}}

\usepackage{fancyhdr}
\pagestyle{fancy}
\fancyhf{}
\renewcommand{\headrulewidth}{0pt}
\fancyfoot[L]{\footnotesize\color{mDim}Vanguard $\cdot$ Meridian Technical Brief $\cdot$ v4 (Stage 6+7 honest)}
\fancyfoot[R]{\footnotesize\color{mDim}\thepage}

\begin{document}
\thispagestyle{empty}

\begin{center}
  {\Huge\bfseries\color{mCyan}Vanguard\,+\,Meridian}\\[3pt]
  {\large\bfseries Autonomous Surface Vessel — Technical Brief}\\[2pt]
  {\color{mDim}\small 19 April 2026 \quad\textbullet\quad Prepared for the Vanguard team \quad\textbullet\quad v4 honest edition}
\end{center}
\vspace{2pt}
\rule{\textwidth}{0.4pt}

% ----- WHERE WE ARE --------------------------------------------------------
\section{What this document is}
A straight accounting of what's built, what's drafted, what's planned for
the first field deploy, and what's the longer-term platform vision. No
overclaims. Status badges: \statusbuilt\ = shipping and tested,
\statustuning\ = shipping with known rough edges, \statusplan\ = on the
work list for first deploy, \statusfuture\ = platform roadmap.

This v4 adds the Stage 6 learned matcher + Stage 7 AIS fusion and
replaces last month's bathymetric results with a larger
cross-resolution evaluation on synthetic and real NOAA ETOPO charts.

% ----- THE GOAL ------------------------------------------------------------
\section{The design goal}
\begin{tcolorbox}[colframe=mCyan, colback=white, title=\normalsize The Vanguard claim — once first deploy is proven]
\textbf{A USV autopilot that uses the Orca Core 2 as its only primary GPS
— no dedicated GPS puck or moving-baseline antenna mast required —
running modern firmware on the Cube Plus the boat already has, and
able to continue navigating when GPS is jammed by fusing depth-chart
matching with AIS-landmark observations.}
\end{tcolorbox}

Every autonomous USV on the market ships with at least one dedicated GPS
puck on the deck, often two for moving-baseline heading. The Orca Core 2
already contains a 10\,Hz GNSS receiver and a 9-axis IMU; if we treat it
as the primary nav source, a whole hardware layer disappears from the
install. When GPS is lost (jam, spoof, canyon), Meridian's bathy +
AIS-landmark filter keeps us navigating.

\pagebreak

% ----- PAGE 2: CURRENT STATE ----------------------------------------------
\section{Current state — what's actually built}
\begin{tabularx}{\textwidth}{@{}l l X@{}}
\toprule
\textbf{Component} & \textbf{Status} & \textbf{Notes} \\
\midrule
\rowcolor{mRow} Tablet GCS (\texttt{mission.html}) & \statusbuilt & AIS overlay, perception, COLREG threat avoid, bench-test, compass-cal wizard, rudder autotune, demo tour, GPS-jam sim, \textbf{bathy marker + AIS-fusion visualization} (dashed cyan lines from estimate to each fusion target). Runs in any browser. \\
\texttt{meridian-orca-bridge} daemon & \statusbuilt & Full package. Mock-mode end-to-end verified. 7/7 unit tests green. \\
\rowcolor{mRow} VESC Rust driver (\texttt{no\_std}) & \statusbuilt & Full COMM\_PACKET, 4/4 tests, integrated in drivers crate (201/201 green). Needs ESC wired UART to exercise on real hardware. \\
Meridian autopilot firmware (24-state EKF, modes, safety, MNP) & \statusbuilt & 74\,k lines Rust, 1{,}200+ tests. SITL-proven across quad / plane / rover / boat / sub. \\
\rowcolor{mRow} Cube Plus board config (\texttt{boards/CubeOrangePlus.toml}) & \statusplan & Drafted. Not compiled, not flashed. First-boot validation is the next major milestone. \\
\texttt{vanguard\_defaults.meridian-params} bundle & \statusbuilt & Matches Meridian naming. Needs live load against running firmware. \\
\rowcolor{mRow} Orca Core 2 wire protocol decoder & \statusplan & YD RAW + Signal K decoders written. Orca-native stubbed — awaiting live capture. \\
Bathymetric map-matching (MCC bootstrap particle filter) & \statusbuilt & \textit{Stage 1--5 complete.} Monte Carlo (N=20 per condition, 4 chart resolutions): see page 3. Deployed as auto-selected filter for charts coarser than 80\,m. \\
\rowcolor{mRow} \textbf{CNN learned feature matcher} (Stage 6) & \statusbuilt & 105\,k-param BatchNorm+Dropout net, 3-channel chart/obs fusion, trained on 308k samples with 10\% sonar-outlier injection. Deployed as primary filter for charts $\leq$25\,m. See page 3 for head-to-head. \\
\textbf{Bathy + AIS-landmark fusion} (Stage 7) & \statustuning & CNN fusion path working: 2 AIS targets drop median error from 6.9\,m $\rightarrow$ 3.9\,m. Bootstrap-path fusion collapses under sharp likelihood; tuning (Cauchy-scale + adaptive regularization) pending. Tablet fusion visualizer ships today. \\
\rowcolor{mRow} Chart loader (NetCDF, GeoTIFF, GEBCO, NOAA ETOPO) & \statusbuilt & Real NOAA ETOPO cache (Botany Bay, 464\,m cells) tested end-to-end. \\
IMU-aided dead reckoning during GPS outage & \statusbuilt & Velocity hold 30\,s + 60\,s exponential decay. \\
\rowcolor{mRow} GPS-jam demo mode on the tablet & \statusbuilt & One-click simulation: GPS goes dark, nav badge flips through GPS $\rightarrow$ BATHY $\rightarrow$ BATHY+AIS$\times$N depending on available landmarks. \\
Documentation bundle (7 docs) & \statusbuilt & Architecture, hardware ID, pitch, competitive, field kit, deployment, bathy tiers. \\
\bottomrule
\end{tabularx}

\section{What's NOT built yet — transparency}
\begin{itemize}
\item Meridian has never been flashed onto a Cube Plus.
\item No integrated Meridian+Orca test on physical hardware — all in SITL / mock / Monte-Carlo.
\item ESC is PWM today; VESC UART is an option, not wired.
\item Bathy filter is an external daemon feeding the tablet. Firmware-side EKF source integration (Tier 3b) is blocked on Cube flash.
\item Bootstrap+AIS fusion has a likelihood-scale issue. CNN+AIS fusion works. Bootstrap-path tuning is 1--2 days of work.
\item \textbf{Zero field hours.} Everything above is in-lab.
\end{itemize}

\pagebreak

% ----- PAGE 3: BATHY NUMBERS ----------------------------------------------
\section{Bathymetric map-matching — honest performance data}

\textbf{Setup:} N=20 trials per condition, 90\,s trajectory each. Particle
filter initialized 30\,m from truth; boat runs 1.5--2.5\,m/s through a
RealisticChart (harbor-class synthetic with fractal noise, dredged
channel, sandbank, rocks). Three conditions per chart: clean (only
Gaussian depth noise $\sigma{=}0.3$\,m), outlier-every-10-steps, and
harsh-outlier-every-5-steps ($\pm$3\,m sonar spikes). Convergence
criterion: position error drops below 30\,m. Divergence: sustained
error ${>}80$\,m in the final 20\,s.

\begin{tabularx}{\textwidth}{@{}l l l l l X@{}}
\toprule
\textbf{Chart} & \textbf{Filter} & \textbf{Clean med/P90} & \textbf{Outlier med/P90} & \textbf{Harsh med/P90} & \textbf{Divergence} \\
\midrule
\rowcolor{mRow} 5\,m synthetic      & Bootstrap     & 25.7 / 31.3\,m & 26.1 / 34.8\,m & 27.0 / 43.1\,m & 0 / 60 \\
5\,m synthetic                      & CNN v3        & 23.4 / 28.9\,m & 22.9 / \textbf{25.7\,m} & 21.2 / \textbf{28.8\,m} & 4 / 60 \\
\rowcolor{mRow} 25\,m synthetic     & Bootstrap     & 25.1 / 30.9\,m & 24.7 / 34.8\,m & 23.3 / 41.7\,m & 0 / 60 \\
25\,m synthetic                     & CNN v3        & 23.1 / 28.9\,m & \textbf{19.4 / 23.9\,m} & 23.7 / \textbf{25.8\,m} & 2 / 60 \\
\rowcolor{mRow} 100\,m synth (GEBCO-class) & Bootstrap & 20.0 / 27.7\,m & — & — & 0 / 20 \\
NOAA ETOPO 464\,m real              & Bootstrap     & 25.3 / 31.6\,m & 25.3 / 31.8\,m & 29.6 / 71.2\,m & 1 / 12 \\
\bottomrule
\end{tabularx}

\vspace{4pt}
\textbf{Read of the numbers.} On fine-resolution charts the CNN matcher
gives a 3--6\,m median improvement and, more importantly, collapses
the p90 tail under outliers from 35--43\,m down to 24--29\,m. That tail
win is the real differentiator — under realistic field noise, the
bootstrap occasionally takes a long single-step excursion that the
learned matcher absorbs.

\textbf{Read of the honesty.} We are not at published SOTA
(RBPF+MCC papers report 3--8\,m on 1--2\,m charts; LSTM-RBPF papers
report 2--5\,m). Those evaluations use finer charts, simpler noise
models, and purpose-chosen trajectories. Our 20--25\,m median across a
440$\times$ chart-resolution range, with zero catastrophic divergences
at the bootstrap-filter level, is what's real today.

\subsection{Stage 7 — AIS landmark fusion}
When GPS is jammed, we can see AIS targets on VHF (passive receive) and
measure range/bearing via radar. Each observation adds a Cauchy-robust
log-likelihood to every particle. Demo result with CNN filter,
2 stationary AIS targets at 400\,m and 600\,m offsets:
\begin{itemize}
\item No AIS: 6.9\,m median
\item +1 AIS target: 6.9\,m (single range is ambiguous around a circle)
\item +2 AIS targets: \textbf{3.9\,m median} ($\sim$2$\times$ improvement)
\end{itemize}
Bootstrap-path fusion still crushes the posterior to a single particle
under sharp range likelihoods; tuning is open work (expected
completion ${\sim}2$ days).

\subsection{Deployment posture}
\texttt{pick\_filter()} auto-selects: CNN for charts $\leq$25\,m,
bootstrap for $\geq$80\,m, either for 25--80\,m with tie-breaker on
recent filter health. Both filters run in parallel on the bridge
daemon; a health watchdog hot-swaps if one diverges.

\textbf{Runtime budget:} bootstrap 5--9\,ms per 1\,Hz step (fits the
Orca CPU at 10\,Hz easily). CNN 140--170\,ms (fits 1\,Hz comfortably;
10\,Hz would need either particle-count reduction or GPU). Calibration
metric (fraction of steps where true error $\leq 1.5\sigma$ reported):
$\approx 1.00$ for both, i.e., the filter reports \emph{conservative}
uncertainty, which is what a downstream EKF wants.

\pagebreak

% ----- PAGE 4: REDUNDANCY, HARDWARE, COMPARISON ---------------------------
\section{Redundancy — built vs.\ planned}
\begin{tabularx}{\textwidth}{@{}l l l X@{}}
\toprule
\textbf{Layer} & \textbf{Status} & \textbf{Holds for} & \textbf{Mechanism} \\
\midrule
\rowcolor{mRow} 1 & \statusbuilt & ${<}\,50$\,ms & Meridian EKF source arbitration — primary $\rightarrow$ serial-puck standby in one estimation cycle. \\
2 & \statusbuilt & ${<}\,60$\,s & IMU + magnetometer + GPS COG fusion gives multi-source heading. \\
\rowcolor{mRow} 3 & \statusplan & ${<}\,5$\,min & VESC motor RPM $\rightarrow$ speed estimate. Driver built, wiring pending. \\
4 & \statusbuilt & no time limit & Bathy map-matching with auto-selected filter. \textbf{20--25\,m median, zero divergences on bootstrap across 60+ trials} per condition. Tablet shows estimate live; Cube-side EKF integration pending Tier 3b. \\
\rowcolor{mRow} 4.5 & \statustuning & no time limit & \textbf{Stage-7 bathy+AIS fusion} — 3.9\,m median with 2 AIS targets (CNN path). Collapses bootstrap-path posterior; tuning in flight. \\
\bottomrule
\end{tabularx}

\section{Where this sits vs.\ the alternatives}
\begin{tabularx}{\textwidth}{@{}L{38mm} l l l X@{}}
\toprule
\textbf{} & \textbf{Install} & \textbf{GPS rate} & \textbf{GPS-denied nav} & \textbf{Key limitation} \\
\midrule
\rowcolor{mRow} Industry dual-GPS USV & 4--6\,h & 5\,Hz & IMU DR only & Moving-baseline rig, fragile \\
SeaRobotics USV-2600 & Factory & Proprietary & Not published & Closed stack, \$100K+ \\
\rowcolor{mRow} Clearpath Heron & Factory & 5\,Hz & None stock & ROS-dependent \\
BlueRobotics BoatBox & DIY & 5\,Hz & None & Legacy autopilot \\
\rowcolor{mCyan!10}\textbf{Vanguard + Meridian (target)} & \textbf{${<}1$\,h} & \textbf{10\,Hz} & \textbf{Bathy + AIS fusion} & \textbf{Requires Orca Core 2; first field deploy pending} \\
\bottomrule
\end{tabularx}

\section{Live demo scenario (for SOCOM)}
One tablet button simulates a GPS jam. Within 2\,s the status bar flips
from \textit{NAV: GPS} (green) to \textit{NAV: BATHY (GPS DENIED)}
(amber). If two or more AIS contacts are within 1.5\,km, the badge
upgrades to \textit{NAV: BATHY\,+\,AIS$\times$2} (cyan) and dashed cyan
lines appear on the map linking the boat's estimated position to each
fusion AIS target, with the measured range labeled on each line. The
boat continues the mission autonomously with no operator intervention.
60\,s later GPS restores and the badge returns to green.

\section{Path forward — first field deploy}
\begin{enumerate}[leftmargin=*, itemsep=3pt]
\item \textbf{Flash Meridian to the Cube Plus} (2--4 days).
\item \textbf{Load parameter bundle via the tablet} (half-day).
\item \textbf{Bench-test with webcam} (half-day). Includes rudder autotune.
\item \textbf{Orca wire-protocol capture} (1 day).
\item \textbf{Compass calibration on-site} (30\,min).
\item \textbf{Dock test} (half-day).
\item \textbf{Sheltered-water mission} (half-day).
\item \textbf{Full field deploy} (1 day). Scripted GPS-denied segment + AIS-fusion visual.
\end{enumerate}
\textbf{Total:} $\sim$7--10 working days from "we start flashing" to
field-data-in-hand.

\section{Open items / known rough edges}
\begin{itemize}
\item \textbf{Bootstrap+AIS fusion collapse.} The likelihood is too sharp relative to the filter's particle spread; Cauchy softening helps but doesn't fully fix it. Plan: adaptive robust-scale tuned to current particle spread, plus a minimum post-resample regularization floor. 1--2 days.
\item \textbf{CNN occasional divergence (3--5\% on fine charts).} Under investigation — likely a single-step collapse when the CNN log-likelihood happens to peak on a wrong particle. Mitigation in place: parallel bootstrap filter as hot-standby via health watchdog.
\item \textbf{Cube flash not done.} Blocks Tier 3b firmware-side EKF integration.
\item \textbf{No in-water hours.} Everything above is lab / Monte Carlo.
\end{itemize}

\section{The bottom line}
\begin{tcolorbox}[colframe=mSafe, colback=white, title=\normalsize What you're signing up for]
\textbf{Today:} mature SITL-proven autopilot; a tablet GCS with
bathy nav, AIS-fusion visualization, and a scripted GPS-jam demo;
an Orca-ingest daemon (mock-mode verified); a validated VESC driver;
a field-robust bathy map-matcher (bootstrap + CNN) that holds to 20--25\,m
median across chart resolutions from 5\,m to 464\,m with zero divergences
on the bootstrap path; and a Stage-7 AIS-landmark fusion that takes the
CNN-path position error to 3.9\,m with two targets.

\textbf{Next 1--2 weeks:} flash Meridian onto the Cube Plus, tune the
bootstrap+AIS path, wire the Orca feed from live hardware, run bench
and dock tests, take the boat out for its first autonomous mission.

\textbf{3-week horizon (SOCOM demo):} in-water validation done, scripted
GPS-denied + AIS-fusion scenario rehearsed, pre-recorded backup footage
in case the live link flakes on the day.

The architecture is real and tested; the hardware integration is the
last mile and we're on it.
\end{tcolorbox}

\end{document}
